\begin{textsample}{NEG Dim 4 – 2013\_c5 – Score 1.00 – 2013\_c5/enade\_2013\_fono\_b6ccec77c7.txt}  \label{ex:f4_neg_129}
CAPÍTULO 5 DISTRIBUIÇÃO DOS CONCEITOS 5.1 PANORAMA NACIONAL DA DISTRIBUIÇÃO DOS CONCEITOS A Tabela 5.1 apresenta a quantidade e distribuição de cursos de Fonoaudiologia participantes do ENADE/2013, por faixa de conceito e Grande Região.
A diferença entre os cursos tabulados neste capítulo e no capítulo 2 são os cursos sem conceito, em princípio, aqueles sem alunos concluintes que participassem da prova ou que tivessem tão somente um aluno concluinte13.
Observando -se os dados da Tabela 5.1, nota -se que, dos 64 cursos participantes, 24 ( 37,5 \% ) classificaram -se com conceito 3, o valor modal.
Este foi também o conceito modal nas regiões Nordeste, Sudeste e Centro-Oeste ( respectivamente 47,1, 42,3 e 33,3 \% ).
O conceito 4 foi o segundo mais frequente em nível nacional ( 21,9 \%, correspondendo a 14 cursos ) e o conceito 2, o terceiro ( 17,2 \%, correspondendo a 11 cursos ).
Houve, ainda, cinco cursos ( 7,8 \% ) que receberam conceito 5 e outros sete cursos que receberam conceito 1.
Dos 64 cursos de Fonoaudiologia, apenas três ficaram sem conceito ( SC ).
A região Norte participou com cinco cursos ou 7,8 \% do total nacional.
Nesta região, o conceito 2 foi atribuído a dois cursos dos cinco participantes, o que equivale a 40,0 \% do total regional.
Apenas um curso ficou sem conceito ( SC ), enquanto nenhum curso recebeu os conceitos 4 ou 5.
Outros dois cursos ficaram distribuídos entre os conceitos 1 e 3 ( um curso, cada ). 13 Estes cursos com somente um concluinte não foram considerados no capítulo 2.
A região Nordeste participou com 17 cursos ou 26,6 \% do total nacional.
Nesta região, oito cursos, 47,1 \% em termos regionais, obtiveram conceito 3, o conceito modal para a região, como já comentado.
Os conceitos 2 e 4 foram atribuídos a, respectivamente, quatro e dois cursos ( 23,5 \% e 11,8 \% ).
Já os conceitos 1 e 5 foram atribuídos a, respectivamente, dois e um cursos ( 11,8 \% e 5,9 \% ).
Nessa região, nenhum dos cursos ficou sem conceito ( SC ).
Como já comentado, dos 26 cursos participantes da região Sudeste, 11 ( 42,3 \% ) obtiveram conceito 3, o conceito modal.
O conceito 1 foi atribuído a três cursos ( 11,5 \% ) e o conceito 2, a apenas dois ( 7,7 \% ).
O conceito 4 foi atribuído a seis cursos ( 23,1 \% ) e, também, dois outros ( 7,7 \% ) receberam o conceito 5.
Apenas dois cursos, também, ficaram sem conceito ( SC ).
A região Sul também contou com cursos distribuídos em quase todas as faixas de conceitos, a exceção do conceito 1 sem nenhum curso.
A predominância do conceito 4 foi de 38,5 \%, correspondentes a cinco dos 13 cursos participantes na região Sul.
Outros seis cursos receberam igualmente os conceitos 2 e 3 ( três cursos, cada ).
O conceito 5 foi recebido por dois cursos ( 15,4 \% ).
Nenhum dos cursos da região Sul ficou sem conceito ( SC ).
De todos os três cursos participantes na região Centro-Oeste, todos receberam cursos igualmente associados aos conceitos 1, 3 e 4 ( um curso, cada ).
Nesta região nenhum curso ficou sem conceito ( SC ) ou recebeu os conceitos 2 ou 5.
O Gráfico 5.1 apresenta a distribuição do Conceito ENADE segundo Grande Região : poligonais mais à esquerda representam Grandes Regiões com pior distribuição, e poligonais mais à direita, Grandes Regiões com melhores conceitos.
A região Norte ( linha azul ) apresentou os cursos em IES com valores mais baixos, ao passo que a região Sul ( linha roxa ) apresentou a distribuição com valores maiores.
Cursos em IES das regiões Nordeste ( linha verde ), Sudeste ( linha vermelha ) e Centro-Oeste ( linha preta ) apresentam uma situação intermediária. 5.2 CONCEITOS POR CATEGORIA ADMINISTRATIVA E POR GRANDE REGIÃO A Tabela 5.2 apresenta a distribuição dos cursos participantes do ENADE/2013 de Fonoaudiologia, por Categoria Administrativa, de acordo com os conceitos por eles alcançados, segundo as Grandes Regiões brasileiras.
Dos 64 cursos participantes, 19 ( 29,7 \% ) eram ministrados em Instituições Públicas e 45 ( 70,3 \% ), em Privadas.
De acordo com as informações da Tabela 5.2, em termos nacionais, das cinco IES que receberam o conceito 5, todas eram públicas.
Dos 19 cursos participantes de IES Públicas, o conceito 4 foi o valor modal, atribuído a sete cursos ( 36,8 \% ).
Dois cursos receberam igualmente os conceitos 1 e 2 ( um curso, cada ).
Entre os demais cursos de Instituições Públicas participantes, também, outros dez obtiveram igualmente os conceito 3 e 5 ( cinco cursos, cada ).
Nesta categoria, nenhum dos cursos ficou sem conceito ( SC ).
Na rede privada, o conceito modal foi 3, com 19 cursos dos 45 da categoria.
Entre os demais cursos participantes, seis receberam conceito 1 e dez, conceito 2.
O conceito 4 foi Na análise por região, observa -se que, na região Norte, as Instituições Privadas participaram com cinco cursos ( totalidade regional ), dos quais apenas um ficou sem conceito ( SC ).
O conceito modal para as Instituições Privadas na região foi 2, com dois atribuído a apenas sete cursos, enquanto nenhum curso recebeu o conceito 5.
Nesta Categoria Administrativa, apenas três dos cursos ficaram sem conceito ( SC ).
Na análise por região, observa -se que, na região Norte, as Instituições Privadas participaram com cinco cursos ( totalidade regional ), dos quais apenas um ficou sem conceito ( SC ).
O conceito modal para as Instituições Privadas na região foi 2, com dois cursos, correspondendo a 40,0 \%.
Outros dois cursos receberam igualmente os conceitos 1 e 3 ( um curso, cada ).
Enquanto nenhum curso recebeu os conceitos 4 ou 5 nessa categoria e região.
Como já mencionado, não houve nenhum curso de Instituição PúblicaPública nesta região.
Na região Nordeste, a rede privada concentrou dez dos 17 cursos participantes, equivalentes a 58,8 \% do total da região, a menor proporção entre todas as regiões.
Nesta Categoria Administrativa na região, predominaram os cursos com conceito 3, cinco cursos correspondendo a 50,0 \% dos cursos oferecidos por IES Privadas no Nordeste.
Apenas quatro cursos receberam o conceito 2 ( 40,0 \% ), enquanto o conceito 1 participou com somente um curso ( 10,0 \% ).
Nenhum curso ficou sem conceito ( SC ) ou recebeu conceitos 4 ou 5 nessa categoria e região.
As Instituições Públicas dessa região participaram com sete cursos ( 41,2 \% ), dos quais três obtiveram conceito 3, o valor modal, seguidos de dois cursos que receberam conceito 4 e outros dois cursos obtiveram igualmente os conceitos 1 e 5 ( um curso, cada ).
Nenhum curso ficou sem conceito ( SC ) ou recebeu o conceito 2 nesta combinação de região e categoria.
Na região Sudeste, a proporção de cursos da rede privada, 73,1 \%, correspondendo a 19 dos 26 cursos participantes.
Nesta categoria, na região Sudeste, o conceito modal foi 3 ( dez cursos ).
Nesta combinação de Categoria Administrativa e Grande Região, apenas dois dos cursos não receberam conceito ( SC ).
Os demais foram avaliados com conceito 1 ( três cursos ) e conceitos 2 e 4 ( dois cursos, cada ).
Nenhum curso recebeu o conceito 5.
Entre os sete cursos em Instituições Públicas na região Sudeste, a categoria modal foi a que recebeu conceito 4 ( quatro cursos ).
Os demais foram avaliados com os conceitos 3 ( um curso ) e 5 ( dois cursos ).
Nenhum curso ficou sem conceito ( SC ) ou recebeu os conceitos 1 ou 2, nessa categoria da região Sudeste.
As Instituições Privadas concentraram oito dos 13 cursos participantes da região Sul, 61,5 \% do total regional.
Desses, quatro ficaram com conceito 4, o conceito modal.
Nesta combinação de Categoria Administrativa e Grande Região, nenhum curso ficou sem receber conceito ( SC ) ou recebeu os conceitos 1 ou 5.
Outros quatro cursos receberam igualmente os conceitos 2 e 3 ( dois cursos, cada ).
As Instituições Públicas na região Sul participaram com cinco cursos ( 38,5 \% ), dos quais nenhum ficou sem conceito ( SC ) ou recebeu o conceito 1.
O conceito modal foi 5, com dois cursos.
Além disso, outros três cursos receberam igualmente os conceitos 2, 3 e 4 ( um curso, cada ) nessa categoria e região.
Na região Centro-Oeste, todos os três cursos participantes eram de Instituições Privadas ( tal como na região Norte ).
Todos os três cursos receberam igualmente os conceitos 1, 3 e 4 ( um curso, cada ).
Nenhum curso ficou sem conceito ( SC ) ou recebeu os

% matched lemmas: 
\end{textsample}
